\documentclass[runningheads]{llncs}

\usepackage[T1]{fontenc}
\usepackage{graphicx}
\usepackage{booktabs}
\usepackage{amsmath}
\usepackage{amssymb}
\usepackage{url}
\usepackage{enumitem}
\usepackage{xcolor}
\usepackage{array}
\usepackage{ragged2e}

\title{
% Conceptual Modeling with Large Language Models in Structured, Semi-structured, and Unstructured Domains 
% SWL: Need to consider using Formal instead of Structured
% NEED TO REVISIT I THINK. 
%
%Augmenting Conceptual Modeling: Exploring the Versatility of Large Language Models in Structured and Unstructured Domains
%
%Formalization Gradient: Empirical Effects of Representation Formality on Local LLM Performance
% Formalization Matters: An Empirical Study of LLMs for Conceptual Modeling Tasks
Formalization Matters: Conceptual Modeling with Large Language Models in High, Semi and Low-Formal Tasks
}

\titlerunning{Formalization Matters}

%  

\author{Anonymous Author(s)}

\authorrunning{Anonymous Author(s)}

\institute{Anonymous Institution(s)\\
\email{anonymous@domain.com}}


\begin{document}

\maketitle

% ══════════════════════════════════════════════════════════════
\begin{abstract}
%Conceptual modeling remains a highly creative and important task for properly representing a domain for which an information system will be developed.  Large Language Models (LLMs) demonstrate remarkable versatility across a wide range of domains. Therefore, it is reasonable to expect that many researchers are investigating their use in conceptual modeling.  Yet the performance of LLMs varies substantially depending on the degree of formalization and structure inherent in the domain. This research analyzes the capabilities of open-weight LLMs in well-formalized domains such as SQL queries, where problems have clearly defined structures and objective solutions.  It compares their performance in less formalized domains, such as law and management, which rely more heavily on interpretation, contextual reasoning, and social norms.  We develop standardized benchmark tasks  for structured, semi-structured, and unstructured domains. We use these along with gold-standard outputs generated or curated using structured protocols and evaluated against automated overlap metrics to assess the quality, consistency, variability, and robustness of LLM-generated results.  Our findings show that, although LLMs achieve substantially higher recall-based scores in domains with explicit symbolic representations, their outputs in less formalized domains achieve lower recall-based quality and are more susceptible to context drift and normative ambiguity. The study shows how domain formalization influences LLM use for modeling or related tasks.

Conceptual modeling tasks vary in how explicitly they constrain valid interpretations, from executable schemas to informal stakeholder scenarios. This paper studies how task formalization affects the quality and consistency of LLM outputs. We evaluate two open-weight LLMs on three conceptual modeling-related tasks: SQL generation over a relational schema, legal clause interpretation using CUAD-derived categories, and management decision-making from incident scenarios. Using gold-standard outputs and recall-based overlap metrics, we compare output quality, variance, consistency, and robustness across high-, semi-, and low-formalization domains. Results show that output quality improves as formalization increases, but variance follows the opposite pattern: high-formal tasks produce higher variance, while low-formal tasks yield lower but more consistently moderate scores. Less formalized tasks are also more susceptible to context drift and normative ambiguity. These findings suggest that formal representations can improve LLM-assisted conceptual modeling and reasoning, but require verification mechanisms rather than assumptions of predictable correctness.


%reasoning  and discusses implications for the deployment of LLMs in decision-making, education, and support systems.
% See other Overleaf for ideas on revising the abstract.
%Conceptual models vary in their degree of formalization, ranging from fully formal schema languages to semi-formal notations and informal natural-language descriptions. We investigate how this formalization gradient affects the quality and consistency of large language models (LLMs) when they reason over representations at different formalization levels. Using two open-weight models (Mistral-7B-Instruct and Llama-3.1-8B-Instruct), we compare performance across three levels: high-formal (SQL generation from a relational schema), semi-formal (legal clause interpretation from structured contract text), and low-formal (management decision-making from natural-language postmortem scenarios). Each task is executed $K\!=\!5$ times to measure output consistency. Our headline finding is a \emph{variance paradox}: output variance \emph{decreases} monotonically as formalization decreases (Levene's $p < 10^{-76}$), contradicting the intuition that less-structured tasks produce more variable outputs. Formal tasks produce bimodal hit-or-miss distributions, while informal tasks yield consistently moderate quality. We also find that output quality decreases significantly with decreasing formalization (Kruskal-Wallis $p < 10^{-256}$ for all metrics), and that response length partially mediates recall-based quality scores. These findings connect the conceptual modeling tradition of formalization levels to empirical LLM behavior, with implications for AI-assisted modeling, schema-driven code generation, and the automation boundary in model-driven engineering.
\end{abstract}

\textbf{Keywords:} conceptual modeling, formalization levels, large language models, model-driven engineering, output consistency, variance paradox

% ══════════════════════════════════════════════════════════════
\section{Introduction}
\label{sec:intro}
\sloppy

Conceptual modeling remains challenging because it requires understanding and representing the application domain \cite{mylopoulos1992conceptual,olive2007conceptual}. Motivated by the broad versatility of LLMs across domains \cite{bommasani2021opportunities,brown2020language}, researchers have begun applying Generative AI to automate the creation of initial conceptual models \cite{camara2023assessment,fill2023conceptual,haerer2023conceptual,DBLP:conf/er/SilvaMCKP25,Storey2025}.  Representations of application domains span a continuum from highly structured, formal representations to loosely structured, informal descriptions \cite{mylopoulos1992conceptual}, although it is not obvious where along this continuum LLM-based automation is reliable.

Generative AI systems perform well on concrete, well-studied tasks such as summarization 
\cite{DBLP:conf/acl/LewisLGGMLSZ20},
%\cite{goyal2022news},
image generation \cite{rombach2022high},
%tutoring \cite{kasneci2023chatgpt},
and code or math benchmarks
\cite{li2022alphacode,DBLP:conf/nips/LewkowyczADDMRS22}.
%\cite{chen2021evaluating,cobbe2021training}.
For richer applications, however, they remain prone to hallucination~\cite{huang2025survey,ji2023survey}, failures of contextual nuance and representational harms~\cite{bender2021dangers}, and uneven performance across reasoning, factuality, and interaction tasks~\cite{bang2023multitask}. The objective of this research is to identify which conceptual modeling tasks LLMs handle well along the structured-to-unstructured continuum.

To carry out our research we build benchmarks for three levels of formalization that capture how structured a domain is. We analyze results for quality and consistency, and discuss the challenges of context drift, variability, and normative ambiguity. We operationalize \emph{formalization} as the degree to which a task's input constraints, solution space, and evaluation criteria are explicitly specified,
%following the conceptual modeling tradition that distinguishes 
based on whether the applications have
formal, semi-formal, and informal representations \cite{mylopoulos1992conceptual,olive2007conceptual,DBLP:journals/isj/WandW93}. SQL generation (formal domain) from a relational schema provides a single correct answer that is verifiable by execution (running a query). Management decisions (low-formal domain) require weighing competing stakeholder interests with no objectively correct solution. Legal interpretations (semi-formal domain) are based on laws but constrained by text and require interpretive reasoning.


%\paragraph{Why formalization should matter.}
Formalization matters for several reasons. 
LLM training corpora are dominated by unstructured prose, but structured content is deliberately \emph{upsampled}: Llama~3 allocates 17\% of tokens to code and 25\% to mathematical and reasoning data, despite having a far smaller share of raw web text \cite{grattafiori2024llama3herdmodels}.
%\cite{brown2020language,grattafiori2024llama3herdmodels}.
This asymmetric exposure may partly explain stronger performance on schema- and code-like tasks (tasks from domains with a high level of formalization).
%This asymmetric exposure plausibly yields deeper per-token internal representations for formal languages than for open-ended prose.
For the task, as the solution space expands, the model must rely less on pattern matching and more on compositional reasoning, precisely where current LLMs are the weakest \cite{dziri2023faith}.

This research makes three main contributions.
%that appear at the intersection of conceptual modeling and AI:
% I'm not sure why you don't like the statement above. --SWL
\begin{enumerate}[noitemsep]
  %\item We operationalize the classical formal/semi-formal/informal domain representation distinction as an experimental variable, instantiated with three task types: SQL generation from a conceptual model of a sales/order-processing database (100 tasks), legal clause interpretation (59~tasks), and management decision-making from retrospective
  % postmortem
  % scenarios.
  \item We describe application domains where modeling is helpful in high, semi, and low levels of formalization for SQL generation, legal clause interpretation, and management decision-making tasks, respectively.
  
  \item We measure output quality and consistency
  %($K\!=\!5$ repeated runs) 
  at all three levels using two open-weight models,
  %(Mistral-7B and Llama-3.1-8B), 
  discovering a variance paradox in which formal tasks produce \emph{higher}, not lower, output variability.

  \item We propose a formalization metric
  %framework
  $F(T)$ along three dimensions and show empirically that formalization level significantly predicts both LLM output quality and variance.
\end{enumerate}
\fussy

% ══════════════════════════════════════════════════════════════
\section{Related Work}
\label{sec:related}

Conceptual modeling research supports a progression from informal domain descriptions to explicit, formally constrained conceptual schemas. Mylopoulos 
%motivates
notes that there are heterogeneous domains in conceptual modeling~\cite{mylopoulos1992conceptual}. Wand and Weber’s ontological analysis of modeling grammars provides a basis for evaluating the representational adequacy of conceptual modeling languages~\cite{DBLP:journals/isj/WandW93}. Olivé considers schemas to be logic-based specifications of domain knowledge, constraints, and derivation rules~\cite{olive2007conceptual}.
%The formalization spectrum introduced by Mylopoulos 
% Steve, what about saying that we considered the notion of formalism as mentioned by Mylopoulos (I haven't checked what he said) to determine how we could distintinguish the tasks based on their fromal nature. ?
% I don't think JM gets us there. --SWL
%\cite{mylopoulos1992conceptual} was grounded in ontological theory by Wand and Weber \cite{DBLP:journals/isj/WandW93} and Oliv\'{e} \cite{olive2007conceptual}:
The representational effectiveness of a conceptual model is influenced by how formal the description of a domain is. This determines what inferences can be validly drawn from a model, and the correctness guarantees available to any reasoning agent operating over it. This suggests that any reasoning system, logic-based or statistical, should perform better over more formalized representations, because formalization constrains the space of valid interpretations.

We consider LLMs for modeling. A key assumption is that LLMs acquire internal representations of the structure regularities of formalized domains, not merely memorized patterns. The most compelling evidence comes from mathematics, where LLMs have moved beyond benchmarks to what some consider to be genuine discovery: FunSearch found new solutions to the cap-set problem \cite{romeraparedes2024funsearch}; AlphaProof, coupling a language model with reinforcement learning over Lean, reached silver-medal level at the 2024 International Mathematical Olympiad \cite{alphaproof2024imo}; and Knuth \cite{knuth2026claude} reports an LLM-assisted construction for Hamiltonian cycle decompositions that he had been unable to find. These successes cluster in highly formalized domains (mathematics,
%\cite{cobbe2021training},
code generation,
%\cite{chen2021evaluating},
and formal proof), whereas open-ended tasks such as compositional generalization degrade sharply \cite{dziri2023faith}. We term this asymmetry the \emph{formalization hypothesis}: LLM output quality degrades monotonically as task formalization decreases.

Many benchmarks evaluate LLMs on tasks with well-defined answers, such as competitive-programming problems evaluated by program behavior \cite{li2022alphacode} and reading-comprehension questions with annotated answer spans \cite{rajpurkar2016squad}. Ouyang et al. \cite{ouyang2023llm} show that LLM outputs can vary across runs even with identical prompts. The relationship between task structure and output variability, however, remains unexplored. We address both gaps by treating formalization as a continuous dimension and measuring quality and variance at three points along it.

% ══════════════════════════════════════════════════════════════
\section{Method}
\label{sec:method}

To explore the effectiveness of using LLMs for modeling, we use multiple open-weight models for modeling tasks with different levels of formalization. We measure the results against benchmarks. Three hypotheses test the effectiveness: 

\begin{description}[noitemsep,leftmargin=2em,style=nextline]
  \item[H1: Quality gradient]
  \textit{LLM output quality decreases
  %monotonically 
  as domain formalization decreases.} Models trained on corpora 
  %that oversample 
  of structured content should produce higher-quality outputs when the solution space is tightly constrained.

  \item[H2: Variance gradient]
  \textit{Output variance across repeated runs increases as domain formalization decreases.} Less-constrained tasks have more degrees of freedom, so successive runs should diverge more.

  \item[H3: Length--quality mediation]
  \textit{Greater response length corresponds to higher recall scores regardless of formalization level.} Longer outputs raise the probability of coincidentally matching gold-standard terms; controlling for length should therefore narrow apparent quality differences between models.
\end{description}


\subsection{Experiments}

To test H1--H3 we require tasks that differ systematically in how closely their inputs constrain acceptable outputs, while holding constant whether tasks are self-contained, have gold standards, and span graded levels of complexity. A task is self-contained when the prompt includes all information needed to produce an acceptable answer, so that the model need not rely on external context beyond the input. We selected one domain task for each level of formalization as follows:

\begin{description}[style=nextline]
  \item[High-formal: SQL generation]
    Generate the correct SQL query given a relational schema derived from a conceptual model of a sales/order database and a natural-language question. The schema generated from the conceptual model has precise syntax and semantics; correctness is verifiable by execution.
    % \cite{li2022alphacode,olive2007conceptual}. 
    This consisted of 100 tasks over the Northwind schema, spanning simple joins to complex multi-table queries with aggregations and subqueries.
    % Do we need to publish the dataset somewhere?

  \item[Semi-formal: Legal clause interpretation]
    For a given legal scenario, determine whether an action is permitted, required, or prohibited under a contractual clause. The clause text constrains the answer but requires interpretive reasoning about the scope and exceptions.
    % \cite{mylopoulos1992conceptual}.
    This consisted of 59 tasks across 11 clause categories (assignment, termination, liability, etc.) drawn from the Contract Understanding Atticus Dataset (CUAD)~\cite{hendrycks2021cuad}.  The conceptual model is semi-formal, with well-defined categories corresponding to entities and natural-language descriptions of relationships and constraints.

  \item[Low-formal: Management decision-making]
    Given a postmortem description of a software incident, recommend a management response when there are multiple stakeholders. No single correct answer exists; quality is assessed by coverage of relevant considerations.
    %\cite{mylopoulos1992conceptual}.  Not sure this fits. --SWL
    This dataset has 31 tasks across 9 incident categories derived from publicly documented reports (e.g., GitHub's 2018 split-brain outage; Knight Capital's 2012 trading loss); gold standards enumerate key considerations without prescribing a unique solution.  The conceptual model is low-formal, being represented mostly in natural language (except for the list of stakeholders).
    % Do we need to reference the dataset?
\end{description}

\noindent Table~\ref{tab:task-structure} summarizes the input structure, prompt structure, and expected output.
All tasks adhere to three principles ensuring valid cross-level comparison: (i) self-contained inputs (Markov property): the input is a sufficient statistic for the correct output, so that $P(\text{correct output} \mid \text{input}, \text{external context}) = P(\text{correct output} \mid \text{input})$ and the model need not condition on any context beyond the prompt;
%\footnote{Equivalently, each task is a horizon-1 decision in which the prompt is the state and the response is the action; tasks that require external knowledge would correspond to partially observable (POMDP) settings in which the model must infer or hallucinate missing context.}
(ii) gold standards for every task\footnote{High-formal gold standards are reference SQL queries generated with ChatGPT 5.1 (OpenAI) tool assistance but fully tested, evaluated, and revised by a human expert. Semi- and low-formal gold standards were generated by Claude Sonnet 4.6 (Anthropic) using structured annotation prompts that enforce self-containment and traceability to the input text. Neither ChatGPT nor Claude is a model under evaluation. Expert validation is in progress; annotation spreadsheets and the scoring protocol are available.}; and (iii) graded complexity, with tasks within each level spanning a complexity range (e.g., simple vs.\ complex queries; moderate vs.\ complex management scenarios).

\begin{table*}[t]
\centering
\caption{Task structure by formalization level.}
\label{tab:task-structure}
%\vspace{-3mm}
\scriptsize
\begin{tabular}{
  >{\RaggedRight\arraybackslash}p{1.0cm}
  >{\RaggedRight\arraybackslash}p{4.2cm}
  >{\RaggedRight\arraybackslash}p{4.0cm}
  >{\RaggedRight\arraybackslash}p{2.5cm}
}
\toprule
Level & Input Structure & Prompt Structure & Expected Output \\
\midrule
High-formal &
  \texttt{schema} (Northwind DDL, shared across all 100 tasks), \texttt{question} (natural-language query) &
  \texttt{Schema: \{schema\}} $\rightarrow$ \texttt{Question: \{question\}} $\rightarrow$ \texttt{SQL:} &
  Raw SQL query (no markdown, no explanation) \\
\addlinespace
Semi-formal &
  \texttt{clause\_text} (verbatim contract clause from CUAD), \texttt{scenario} (factual situation), \texttt{question} (legal interpretation question) &
  \texttt{Contract Clause: \{clause\_text\}} $\rightarrow$ \texttt{Scenario: \{scenario\}} $\rightarrow$ \texttt{Question: \{question\}} $\rightarrow$ \texttt{Answer:} &
  Legal interpretation citing specific clause terms \\
\addlinespace
Low-formal &
  \texttt{scenario} (postmortem incident), \texttt{stakeholders} (named roles), \texttt{question} (management decision question) &
  \texttt{Scenario: \{scenario\}} $\rightarrow$ \texttt{Key Stakeholders: \{stakeholders\}} $\rightarrow$ \texttt{Question: \{question\}} $\rightarrow$ \texttt{Analysis and Recommendation:} &
  Multi-perspective analysis addressing stakeholder interests and tradeoffs \\
\bottomrule
\end{tabular}
\end{table*}


%\subsection{Formalization Metric Framework}
%\label{sec:formalization-metric}

%For experimental purposes,
To empirically test our hypotheses,
we combine the three formalization levels into a continuous \textit{formalization metric} $F(T)$ along three dimensions. \textit{Output-space cardinality $C(T)$} is the number of distinct outputs judged correct: small for SQL (one canonical query plus syntactic variants), bounded for legal interpretation, and effectively unbounded for management decisions. \textit{Syntactic constraint density $S(T)$} is the proportion of the output determined by grammar rules of the target language rather than by model choice (high for SQL, loose rhetorical conventions for legal, none beyond natural language for management). \textit{Evaluation decidability $D(T)$} captures whether correctness can be determined algorithmically ($D{=}1$, SQL execution), by bounded expert judgment ($D{=}0.5$, legal clause matching), or only by open-ended assessment ($D{=}0$, management quality).

We define $F(T) = \tfrac{1}{3}(\hat{C}(T) + S(T) + D(T))$, where $\hat{C}(T) = 1 / \log_2(1 + C(T))$ normalizes cardinality to $[0,1]$. The denominator has an information-theoretic reading: $\log_2 C(T)$ is the Shannon entropy of a uniform distribution over the $C(T)$ acceptable outputs. The $+1$ offset makes the unique-answer case ($C{=}1$) map to $\hat{C}{=}1$ rather than a singularity, while the logarithm ensures diminishing returns. Thus, moving from one to two defensible answers reduces formalization more than moving from a hundred to two hundred.

Crucially, $F(T)$ is computed \emph{per task instance}, not merely per level. For each of the 190 tasks, $C(T)$, $S(T)$, and $D(T)$ are derived programmatically from task metadata and gold-output structural features (e.g., $S(T)$ from SQL-keyword and punctuation density on the gold query for high-formal tasks, and from analogous markers for the other levels), yielding a continuous $F$ value that varies \emph{within} each level as well as across levels. Table~\ref{tab:F-dims} reports the per-dimension values, and Section~\ref{sec:F-validation} tests whether this instance-level $F(T)$ carries explanatory power beyond the categorical high/semi/low encoding.

\begin{table}[ht]
\centering
\caption{Formalization metric $F(T)$ by level: per-dimension inputs and resulting $F$ (mean $\pm$ std over task instances). $\hat{C}$ is normalized output-space cardinality, $S$ syntactic constraint density, $D$ evaluation decidability. $F$ varies within levels because $C(T)$ and $S(T)$ are computed per instance.}
\label{tab:F-dims}
\setlength{\tabcolsep}{6pt}
\begin{tabular}{lrrrrl}
\toprule
Level & $n$ & $\hat{C}$ & $S$ (mean) & $D$ & $F(T)$ \\
\midrule
High & 100 & 0.558 & 0.880 & 1.0 & $0.813 \pm 0.078$ \\
Semi &  59 & 0.273 & 0.423 & 0.5 & $0.399 \pm 0.029$ \\
Low  &  31 & 0.141 & 0.255 & 0.0 & $0.132 \pm 0.008$ \\
\bottomrule
\end{tabular}
\end{table}

The experiments performed for this study of LLM performance across various modeling tasks run locally using open-weight models.  We use two open-weight instruction-tuned models: Mistral-7B-Instruct-v0.3 (7.24B parameters, sliding-window and grouped-query attention)~\cite{jiang2023mistral} and Llama-3.1-8B-Instruct (8.03B parameters, RoPE, 128K context window)~\cite{grattafiori2024llama3herdmodels}. Both run locally on an HTCondor-managed HPC cluster at Anonymous University. Each job uses a Docker container (\texttt{nvcr.io/nvidia/pytorch:25.01-py3}) with one NVIDIA A100 40GB GPU, 4 CPUs, 32GB RAM, and 50GB disk. All inference is performed locally; no API calls to external services are used.

All inference flows through a unified \texttt{LocalChatModel} class that handles model loading, prompt formatting, and generation. The wrapper:

\begin{enumerate}[noitemsep]
  \item Loads the model in FP16 with HuggingFace \texttt{AutoModelForCausalLM}.
  \item Applies the model's chat template for instruction formatting.
  \item Strips the input prompt from the generated output to isolate the model's response.
\end{enumerate}

To assess output consistency, each task is executed $K\!=\!5$ times with identical prompts, temperature ($T\!=\!0.7$), and a 2048-token cap. Results are stored in JSONL with a \texttt{run\_index} field (0--4) for consistency analysis. A bootstrap power analysis confirms that $K\!=\!5$ is sufficient: even $K\!=\!1$ yields significance $99.6\%$ of the time for the cross-level Kruskal-Wallis test. Table~\ref{tab:dataset} summarizes the experimental data.

\begin{table}[ht]
\centering
\caption{Dataset summary across formalization levels and models.}
\label{tab:dataset}
\setlength{\tabcolsep}{5pt}
\begin{tabular}{llrrrr}
\toprule
Level & Model & Tasks & $K$ & Records & Avg. Resp. Len. \\
\midrule
High-formal & Llama-3.1-8B   & 100 & 5 & 500 & 967 \\
High-formal & Mistral-7B     & 100 & 5 & 500 & 212 \\
Semi-formal & Llama-3.1-8B   &  59 & 5 & 295 & 2357 \\
Semi-formal & Mistral-7B     &  59 & 5 & 295 & 833 \\
Low-formal  & Llama-3.1-8B   &  31 & 5 & 155 & 2908 \\
Low-formal  & Mistral-7B     &  31 & 5 & 155 & 2557 \\
\midrule
\multicolumn{2}{l}{Total} & 190 & & 1900 & \\
\bottomrule
\end{tabular}
\end{table}


\subsection{Evaluation}
\label{sec:metrics}

To evaluate our experimental results, we used quality and consistency metrics, applying appropriate statistical tests to assess significance.  We also developed a rubric for human expert evaluation to cross-check the automated metrics.

%rather than metric incommensurability.

\vspace{5pt}
\noindent
\textit{Quality metrics.}
A challenge is finding metrics applicable uniformly across all three levels. Domain-specific metrics (execution correctness for SQL, inter-anno\-ta\-tor agreement for legal) cannot be compared across levels. Therefore we adopt text-overlap metrics against expert gold standards. While coarser than domain-specific alternatives, this ensures that observed differences reflect the formalization gradient rather than incompatibilities among metrics. We use three recall-oriented overlap measures against the gold standard, where $W$ denotes the set of unique lowercased words---$W_{\text{gold}}$ for words from the gold standard and $W_{\text{pred}}$ for words from the LLM prediction (output):
\begin{itemize}[itemsep=4pt]
  \item \textit{Word overlap}: $|W_{\text{pred}} \cap W_{\text{gold}}| / |W_{\text{gold}}|$, capturing broad topical coverage.
  \item \textit{Bigram overlap}: same ratio over word bigrams, sensitive to phrasal structure (e.g., \texttt{GROUP BY} vs.\ separate \texttt{GROUP} and \texttt{BY}; ``material breach'' vs.\ unrelated occurrences of either word).
  \item \textit{Key-term recall}: fraction of distinctive gold-standard words (length $>3$, excluding stopwords) present in the prediction, focusing on domain-specific content words (table names, legal terms, stakeholder labels).
\end{itemize}
All three are recall-biased: they measure how much of the gold standard is recovered, not how much of the prediction is relevant. This is deliberate but introduces a known bias toward verbose responses (see H3 and Section~\ref{sec:limitations}). Recall-based quality metrics by formalization level and model are summarized in Table~\ref{tab:quality}.

\begin{table}[htbp]
\centering
\caption{Quality metrics (recall-oriented overlap against gold standards) by formalization level and model. All levels are computed on a single raw-text word-set representation (letters only, length $\geq 3$; key-terms length $>4$ minus stopwords), resolving a prior representation mismatch in which high-formal used raw output while semi/low used a legacy computation. Pred/Gold is the mean character-length ratio. See \texttt{scripts/recompute\_unified\_metrics.py}. Precision and $F_1$ for these same sets appear in Table~\ref{tab:precision}.}
\label{tab:quality}
\setlength{\tabcolsep}{5pt}
\begin{tabular}{llrrrrr}
\toprule
Level & Model & $N$ & Word$\uparrow$ & Bigram$\uparrow$ & KeyTerm$\uparrow$ & Pred/Gold \\
\midrule
High  & Llama   & 500 & .800 & .510 & .791 & 7.00$\times$ \\
High  & Mistral & 500 & .764 & .442 & .748 & 1.06$\times$ \\
Semi  & Llama   & 295 & .251 & .070 & .181 & 1.73$\times$ \\
Semi  & Mistral & 295 & .181 & .044 & .124 & 0.61$\times$ \\
Low   & Llama   & 155 & .231 & .052 & .204 & 0.64$\times$ \\
Low   & Mistral & 155 & .229 & .053 & .205 & 0.57$\times$ \\
\bottomrule
\end{tabular}
\end{table}

% \vspace{5pt}
\noindent
\textit{Consistency metrics ($K\!=\!5$).}
We measure \textit{unique outputs}, the number of distinct responses across $K$ runs per task after lowercasing and whitespace normalization ($K$ = complete inconsistency; $1$ = perfect reproducibility); and \textit{perfect consistency}, the percentage of tasks producing identical output across all $K$ runs (the strictest reliability criterion).

\vspace{5pt}
\noindent
\textit{Statistical tests.}
Shapiro-Wilk tests reject normality for all quality distributions, so we adopt a non-parametric pipeline operating on ranks. A Kruskal-Wallis $H$ test (omnibus test of H1) checks whether central tendencies differ across the three levels; Mann-Whitney $U$ tests with Bonferroni correction ($k\!=\!3$) localize pairwise effects, with high$\,\to\,$semi and semi$\,\to\,$low as the theoretically motivated transitions; Levene's test, which is relatively robust to departures from normality, addresses H2 by testing equality of spreads; and rank-biserial correlation $r \in [-1,+1]$ quantifies effect size so that statistical significance is distinguished from practical magnitude.

\vspace{5pt}
\noindent
\textit{Human evaluation protocol.}
Because automated overlap metrics are proxies for semantic quality, we designed an expert rubric applicable across all three levels:

\begin{description}[noitemsep]
  \item[3 (Excellent):] Correct and complete: SQL executes and returns the right result set; legal identifies all applicable provisions and reaches the correct conclusion; management covers all key stakeholders and tradeoffs.
  \item[2 (Adequate):] Substantially correct; captures the core answer but misses secondary considerations.
  \item[1 (Partial):] Some relevant content but missing major elements or containing significant errors.
  \item[0 (Incorrect):] Wrong, irrelevant, or non-responsive.
\end{description}

\noindent Each output is scored independently by two annotators; inter-annotator agreement is measured by Cohen's kappa (target $\kappa \geq 0.70$). A stratified sample of 30 outputs (10 per level, balanced across models and complexity tags) is scored in a pilot round to calibrate the rubric before full annotation. Expert evaluation has begun through random checks, but the full rubric-based annotation has not yet been fully completed.
%(Section~\ref{sec:human-results}).

% ══════════════════════════════════════════════════════════════
% ══════════════════════════════════════════════════════════════
\section{Results}
\label{sec:results}

\subsection{Quality and Variance Metrics (H1 and H2)}
\label{sec:quality}

To assess H1, we apply statistical tests and perform pairwise comparisons between the metrics described in Table~\ref{tab:quality}, which shows that quality drops sharply with decreasing formalization. Averaging across models, word overlap falls from $.782$ (high) to $.216$ (semi) and $.230$ (low); the dominant gradient is the high-to-semi drop, after which semi- and low-formal tasks occupy a similar quality band (the small semi--low differences are examined in Section~\ref{sec:model-comparison}). The Kruskal-Wallis test reported in Table~\ref{tab:kruskal} confirms that the across-level gradient is highly significant for all metrics. All overlap and length metrics are computed on a single raw-text representation (Table~\ref{tab:quality}); an earlier version of the analysis mixed raw and SQL-normalized text at the high-formal level, and the unified representation is used throughout.

\begin{table}[htbp]
\centering
\caption{Kruskal-Wallis test for quality differences across formalization levels.}
\label{tab:kruskal}
\setlength{\tabcolsep}{5pt}
\begin{tabular}{lrrl}
\toprule
Metric & $H$ & $p~~~~~~$ & Sig. \\
\midrule
Word overlap    & 1410.32 & $5.7 \times 10^{-307}$ & *** \\
Bigram overlap  & 1313.20 & $6.9 \times 10^{-286}$ & *** \\
Key-term recall & 1450.51 & $< 10^{-300}$ & *** \\
\bottomrule
\multicolumn{4}{l}{\scriptsize Significance codes: *** $p < 0.001$.}
\end{tabular}
\end{table}

%\paragraph{Pairwise comparisons.}
Mann-Whitney $U$ tests with Bonferroni correction reveal that the dominant quality drop occurs between high-formal and semi-formal levels. (See Table~\ref{tab:pairwise}.) The semi$\,\to\,$low transition reveals smaller, metric-dependent differences.

\begin{table}[htbp]
\centering
\caption{Pairwise Mann-Whitney $U$ tests (Bonferroni-corrected, $k\!=\!3$). Effect size $r$ is rank-biserial correlation.}
\label{tab:pairwise}
\setlength{\tabcolsep}{5pt}
\begin{tabular}{llrrrrl}
\toprule
Metric & Comparison & $\Delta$ & $U$ & $p_{\text{corr}}$ & $r~~$ & Sig. \\
\midrule
Word$\uparrow$
  & High$\,\to\,$Semi & $+.566$ & 586827 & $1.9 \times 10^{-238}$ & $-.989$ & *** \\
  & High$\,\to\,$Low  & $+.552$ & 309533 & $4.4 \times 10^{-155}$ & $-.997$ & *** \\
  & Semi$\,\to\,$Low  & $-.014$ & 63238 & $8.0 \times 10^{-14}$ & $+.308$ & *** \\
\midrule
Bigram$\uparrow$
  & High$\,\to\,$Semi & $+.419$ & 573804 & $1.1 \times 10^{-217}$ & $-.945$ & *** \\
  & High$\,\to\,$Low  & $+.424$ & 305482 & $4.8 \times 10^{-147}$ & $-.971$ & *** \\
  & Semi$\,\to\,$Low  & $+.005$ & 53534 & $4.3 \times 10^{-24}$ & $+.415$ & *** \\
\midrule
KeyTerm$\uparrow$
  & High$\,\to\,$Semi & $+.618$ & 588186 & $1.0 \times 10^{-240}$ & $-.994$ & *** \\
  & High$\,\to\,$Low  & $+.565$ & 309747 & $1.4 \times 10^{-155}$ & $-.998$ & *** \\
  & Semi$\,\to\,$Low  & $-.052$ & 40072 & $3.1 \times 10^{-43}$ & $+.562$ & *** \\
\bottomrule
\multicolumn{4}{l}{\scriptsize Significance codes: *** $p < 0.001$.}
\end{tabular}
\end{table}

% \subsection{Output Variance Across Formalization Levels (H2)}
% \label{sec:variance}

To assess H2, we use 
Levene's test. Table~\ref{tab:levene} shows a striking finding: output variance \emph{decreases} monotonically with decreasing formalization.
The IQR (interquartile range) follows the same pattern: $.198/.110/.041$ (word overlap) from high to low. This suggests that high-formal tasks produce a bimodal distribution (either close to the gold-standard SQL or completely wrong), while low-formal tasks consistently cover similar conceptual ground even when wording differs.

\begin{table}[ht]
\centering
\caption{Levene's test (median-centered) for equality of variances across formalization levels, on the unified raw-text representation (Table~\ref{tab:quality}). Variance shown per level; the gradient is high $\gg$ semi $\gg$ low for all metrics.}
\label{tab:levene}
\setlength{\tabcolsep}{6pt}
\begin{tabular}{lrrrrrl}
\toprule
Metric & $W$ & $p$ & Var$_{\text{High}}$ & Var$_{\text{Semi}}$ & Var$_{\text{Low}}$ & Sig. \\
\midrule
Word overlap    & 186.39 & $1.2 \times 10^{-74}$  & .0226 & .0115 & .0008 & *** \\
Bigram overlap  & 551.39 & $1.7 \times 10^{-189}$ & .0548 & .0054 & .0002 & *** \\
Key-term recall & 213.97 & $1.6 \times 10^{-84}$  & .0240 & .0118 & .0009 & *** \\
\bottomrule
\multicolumn{4}{l}{\scriptsize Significance codes: *** $p < 0.001$.}
\end{tabular}
\end{table}

\paragraph{Separating across-run from across-task variance.}
H2 concerns variance \emph{across repeated runs}, but the pooled per-record variance in Table~\ref{tab:levene} also absorbs \emph{between-task} dispersion (how much tasks differ from one another within a level). These are distinct quantities, and only the former is what H2 names. We therefore decompose them (Table~\ref{tab:variance-decomp}): for each task we compute the variance of word overlap across its $K\!=\!5$ runs (within-task, across-run), and separately the variance of per-task means (between-task). The pooled variance is dominated by the between-task component (e.g., at the high-formal level $.0187$ of the $.0226$ pooled variance is between-task), so the original Levene test primarily reflected task heterogeneity rather than run-to-run instability. Critically, however, the across-run variance---the quantity H2 actually predicts---exhibits the \emph{same} monotonic decrease with formalization ($.0050 \to .0017 \to .0005$; Kruskal-Wallis $H\!=\!51.1$, $p\!=\!7.9\times10^{-12}$). The refutation of H2 is thus robust to the measurement correction: whether measured within tasks or between tasks, variance decreases as formalization decreases.

\begin{table}[ht]
\centering
\caption{Variance decomposition of word overlap by level (pooled over models). \emph{Within-task} is the mean over tasks of the variance across the $K\!=\!5$ runs of a task (the quantity H2 names); \emph{between-task} is the variance of per-task means; \emph{pooled} is the variance over all run-level records (as used in Table~\ref{tab:levene}). All three components decrease monotonically with formalization.}
\label{tab:variance-decomp}
\setlength{\tabcolsep}{7pt}
\begin{tabular}{lrrrr}
\toprule
Level & \#Tasks & Within-task & Between-task & Pooled \\
\midrule
High & 100 & .00502 & .01866 & .02260 \\
Semi &  59 & .00174 & .01013 & .01145 \\
Low  &  31 & .00045 & .00045 & .00080 \\
\bottomrule
\multicolumn{5}{l}{\scriptsize Kruskal-Wallis on within-task variance across levels: $H\!=\!51.1$, $p\!=\!7.9\times10^{-12}$.}
\end{tabular}
\end{table}

\subsection{Consistency and Length (H3)}
\label{sec:consistency}

Neither Llama nor Mistral achieves meaningful consistency across runs. Averaging across levels, nearly every task produces $K$ unique outputs across $K\!=\!5$ runs (average unique outputs $\approx 5.0$ for all conditions except Mistral on high-formal SQL, where the average is 3.8). The only exception is Mistral-7B on high-formal SQL, where 7 out of 100 tasks (7\%) produce identical outputs across all runs. Perfect consistency is 0\% for all other level/model combinations.

% \subsection{Model Comparison and Length--Quality Mediation (H3)}
\label{sec:model-comparison}

On the recall-oriented metrics, Llama-3.1-8B scores higher than Mistral-7B at the high- and semi-formal levels (Mann-Whitney $U$, Bonferroni-corrected $p < 10^{-5}$), and the two models converge at the low-formal level ($p > 0.99$). This recall advantage, however, is an artifact of verbosity, and it reverses under precision (Table~\ref{tab:precision}). Because the recall metrics divide by $|W_{\text{gold}}|$ and do not penalize irrelevant prediction content, Llama's much longer outputs (high-formal: $7.0\times$ gold length vs.\ Mistral's $1.1\times$) inflate its recall. Precision---$|W_{\text{pred}} \cap W_{\text{gold}}| / |W_{\text{pred}}|$---and the harmonic-mean $F_1$ tell the opposite story at the high-formal level, where Mistral's concise queries achieve far higher precision ($.823$ vs.\ $.547$ word precision) and a higher $F_1$ ($.777$ vs.\ $.615$). Under $F_1$, then, Mistral is the stronger high-formal model, not Llama. At the semi-formal level the models are near-tied under $F_1$ ($.225$ vs.\ $.215$), and at the low-formal level they remain statistically indistinguishable on every measure. This shows that the apparent inter-model quality gap is largely a length effect rather than a difference in task competence.

\begin{table}[htbp]
\centering
\caption{Recall, precision, and $F_1$ by level and model, on the same unified raw-text word-set representation as Table~\ref{tab:quality} (the recall column matches Table~\ref{tab:quality}). Llama's recall lead at the high-formal level reverses under precision and $F_1$: its longer outputs recover more gold terms (recall) but at much lower precision.}
\label{tab:precision}
\setlength{\tabcolsep}{5pt}
\begin{tabular}{llrrrrrr}
\toprule
& & \multicolumn{3}{c}{Word} & \multicolumn{3}{c}{Key-term} \\
\cmidrule(lr){3-5}\cmidrule(lr){6-8}
Level & Model & Rec. & Prec. & $F_1$ & Rec. & Prec. & $F_1$ \\
\midrule
High & Llama   & .800 & .547 & .615 & .791 & .555 & .620 \\
High & Mistral & .764 & .823 & \textbf{.777} & .748 & .813 & \textbf{.763} \\
Semi & Llama   & .251 & .212 & .225 & .181 & .153 & .162 \\
Semi & Mistral & .181 & .273 & .215 & .124 & .195 & .149 \\
Low  & Llama   & .231 & .398 & .291 & .204 & .352 & .257 \\
Low  & Mistral & .229 & .396 & .290 & .205 & .353 & .259 \\
\bottomrule
\end{tabular}
\end{table}

With precision now in hand, H3 (length--quality mediation) is confirmed in its mechanism: greater response length raises recall-oriented scores without a corresponding gain in precision. The inter-model recall gap is largest where the length differential is largest (high-formal) and vanishes where both models produce comparably short outputs (low-formal, $0.6\times$ for both). Crucially, once irrelevant content is penalized, the verbosity premium not only disappears but inverts at the high-formal level, so Llama's higher recall reflects lexical overlap inflation rather than superior understanding. This resolves the qualification noted in earlier analyses that H3 could not be confirmed without precision-sensitive metrics.

\subsection{Robustness and Validity Checks}
\label{sec:within-domain}

A potential confound in the cross-level comparison is that the three formalization levels use different task domains (databases, law, management), so observed quality differences might reflect domain difficulty rather than formalization \textit{per se}. To address this, we exploit the complexity tags embedded in our task sets to test whether the formalization gradient operates \emph{within} a single domain.

\vspace{5pt}
\noindent
\textit{High-formal (SQL).}
The 100 SQL tasks are tagged by difficulty: easy (20 tasks), medium (55 tasks), and hard (25 tasks). These difficulty levels represent a within-domain formalization gradient: easy queries have a more constrained output space ($C(T)$ closer to 1) than hard multi-table queries where multiple valid join strategies exist. Comparing word overlap across difficulty tags, we observe a decline from $.831$ (easy) to $.798$ (medium) to $.708$ (hard), mirroring the cross-level pattern within a single domain.

\vspace{5pt}
\noindent
\textit{Low-formal (management).}
The 31 management tasks span 9 incident categories with two complexity levels (moderate and complex). Complex scenarios, which involve more stakeholders and competing objectives, yield lower gold-standard overlap ($.223$ vs.\ $.240$ for moderate), consistent with the prediction that expanding the solution space reduces overlap-based quality.

These within-domain results strengthen the claim that formalization, not domain, drives the quality gradient. The effect size is smaller within domains than across domains, as expected, since the within-domain formalization range is narrower.

% \subsection{Human Evaluation}
% \label{sec:human-results}

% The unified 0--3 expert rubric described in Section~\ref{sec:metrics} has not yet been fully applied to the experimental outputs, though spot checks have been done. Specifically, the plan is that a sample of 30 outputs (10 per formalization level, balanced across models and complexity tags) will be scored independently by two annotators in a pilot round to calibrate the rubric, followed by full annotation across the complete dataset. Inter-annotator agreement will be assessed via Cohen's kappa (target $\kappa \geq 0.70$).
% %Until expert evaluation is complete,
% The findings we report are currently based on automated overlap metrics. 
% %, whose limitations are discussed in Section~\ref{sec:limitations}.

% \subsection{Robustness of Gold Standards}
% \label{sec:gold-robustness}

A potential threat to construct validity arises when gold standards are generated by the same model family under evaluation, because recall-based overlap metrics may reward shared stylistic tendencies rather than genuine task understanding (i.e., models from the same family often exhibit similar lexical preferences, discourse structures, and reasoning phrasing). In this research, high-formal gold standards are human-curated reference SQL queries, initially created using ChatGPT 5.1 (OpenAI) tool assistance, but then fully tested, evaluated, and revised by a human expert and executed against an actual database. Semi-formal and low-formal gold standards were generated by Claude Sonnet 4.6 (Anthropic), a model architecturally and organizationally independent of both models under evaluation (Mistral-7B and Llama-3.1-8B). This reduces the risk of self-referential contamination: neither evaluated model could benefit from stylistic overlap with its own prior outputs in the gold standard. Additionally, the structured annotation prompts enforce self-containment (all gold-standard content must be traceable to the input text) and explicit reasoning structure, reducing the influence of model-specific stylistic preferences on the gold-standard answers. The quality gradient (H1) is consistent across all three levels, further supporting the robustness of the finding.

\subsection{Validating the Formalization Metric $F(T)$}
\label{sec:F-validation}

A natural objection is that $F(T)$ merely relabels the three levels, adding no information beyond the categorical high/semi/low encoding. Because $F(T)$ is computed per task instance (Section~\ref{sec:method}), we can test this directly with a nested-model ablation over all 190 tasks, regressing per-task mean quality (word overlap) on: M1, the categorical level alone; M2, $F(T)$ alone; and M3, level $+\,F(T)$. If $F(T)$ carried no information beyond the level labels, M3 would not improve on M1.

It does. M3 significantly improves fit over M1 (nested $F_{1,186}=9.31$, $p=2.6\times10^{-3}$; $\Delta\text{AIC}=-7.3$), with a positive within-model coefficient on $F(T)$ ($\beta_F=0.408 \pm 0.134$, $t=3.05$, $p=2.6\times10^{-3}$). Adding a per-level $F$ slope (M4) does not further improve fit ($p=0.77$), so a single global slope suffices. The instance-level $F(T)$ therefore predicts output quality \emph{beyond} the level to which a task belongs. In contrast, $F(T)$ does \emph{not} predict per-task quality \emph{variance} beyond the level encoding ($p=0.65$), consistent with the refutation of H2: variance is governed by level-scale bimodality rather than by fine-grained formalization.

Within-level Pearson correlations between $F(T)$ and mean quality are positive at all three levels (high $r=0.22$, $p=0.03$; semi $r=0.23$, $p=0.08$; low $r=0.42$, $p=0.02$), showing that the signal is not purely a between-level artifact. We further guard against a circularity concern---that both $F(T)$ (through gold-based $S(T)$) and word overlap (through $|W_{\text{gold}}|$) depend on gold-text features---by partialling out gold token count and keyword fraction. The within-level $F(T)$ signal survives robustly at the low-formal level (partial $R^2=0.14$; coefficient inflation $2.2\%$) but is substantially attenuated at the high- and semi-formal levels (inflation $168\%$ and $78\%$; partial $R^2 \leq 0.02$), indicating that the within-level predictive value of $F(T)$ is strongest where the output space is least syntactically constrained. Accordingly, we treat the \emph{cross-level} ablation (M3 vs.\ M1) as the primary evidence that $F(T)$ adds information beyond the categorical encoding, and report the within-level correlations with this caveat.

% ══════════════════════════════════════════════════════════════
\section{Discussion}
\label{sec:discussion}

This research has investigated the performance of LLMs in conceptual modeling by applying LLMs to various modeling cases that are related to different modeling applications that are of a low, semi, or high level of formalization. 
%
%Performance of LLMs varies substantially depending on the degree of formalization and structure inherent in the domain. This paper analyzes the capabilities of open-weight LLMs in well-formalized domains such as SQL queries, where problems have clearly defined structures and objective solutions.  It compares their performance in less formalized domains, such as law and management, which rely more heavily on interpretation, contextual reasoning, and social norms.  Using standardized benchmark tasks we  developed for structured and unstructured domains, with gold-standard evaluations developed by experts, we assess both the quality and consistency of LLM-generated results, as well as the models’ reasoning transparency and robustness.  Our findings show that, although LLMs achieve near-human
%
For the conceptual modeling community, our results show that the differences in the formalization level of a modeling application
%\cite{mylopoulos1992conceptual}
is not simply a matter of convenience. Rather, formalization can serve as a predictor of both LLM output quality and variability. The formalization metric $F(T)$  provides a way to quantify where a task appears on the formalization gradient spectrum and, consequently, what quality and variance 
%profile
to expect from LLM-based creation or analysis of a model.

%PLEASE CHECK IF ABOVE SENTENCE IS OK.
% Looks good to me (I added "spectrum") --SWL

%For the conceptual modeling community, our results establish that the differences in level of formalism \cite{mylopoulos1992conceptual} is not merely a taxonomic convenience but a predictor of both LLM output quality and output variability. The formalization metric framework $F(T)$ proposed in this paper provides a principled basis for quantifying where a task sits on the formalization gradient and, consequently, what quality and variance profile to expect from LLM-based automation.

%\subsection{The Variance Paradox: Formalization Raises the Ceiling but Widens the Distribution}

\subsection{Findings}

The findings as they relate to the hypotheses are summarized in 
%Quality gradient.
%LLM output quality decreases as domain formalization decreases. Models
%trained on corpora of structured content should produce higher-quality out-puts when the solution space is tightly constrained.
%H2: Variance gradient.
%Output variance across repeated runs increases as domain formalization de-
%creases. Less-constrained tasks have more degrees of freedom, so successive
%runs should diverge more.
%H3: Length–quality mediation.
%Greater response length corresponds to higher recall scores regardless of
%formalization level. Longer outputs raise the probability of coincidentally
%matching gold-standard terms; controlling for length should therefore nar-
%row apparent quality differences between models.
Table~\ref{tab:hyp}.

\begin{table}[htbp]
\centering
\caption{Results of testing hypotheses}
\label{tab:hyp}
\setlength{\tabcolsep}{6pt}
\begin{tabular}{c>{\RaggedRight\arraybackslash}p{5.6cm}l}
\toprule
Hypothesis & Description & Support  \\
\midrule
H1 & Output quality decreases as domain formalization decreases & Supported \\
\addlinespace[2pt]
H2 & Output variance increases as domain formalization decreases & Not supported \\
\addlinespace[2pt]
H3 & Greater length has higher recall independent of formalization level & Supported  \\
\bottomrule
\end{tabular}
\end{table}

Of the three hypotheses derived from the formalization hypothesis and related work, two are supported and one is refuted. H1 (quality gradient) is supported: output quality declines sharply as formalization decreases, with the dominant drop at the high-to-semi transition and a near-plateau (with a mild reversal) between semi and low.
%($p < 10^{-256}$).
H3 (length--quality mediation) is supported: the inter-model recall gap covaries with the response-length differential and, once precision is computed (Table~\ref{tab:precision}), the verbosity premium is shown to be lexical rather than substantive---Llama's high-formal recall lead inverts under precision and $F_1$. H2 (variance gradient) is refuted, yielding the most consequential finding: output variance \emph{decreases} rather than increases with decreasing formalization.
%($p < 10^{-76}$). 
Formal tasks produce bimodal 
%hit-or-miss 
distributions, whereas informal tasks yield consistently moderate quality.

% % FROM STEVE: I think we need to discuss this paragraph.  I don't know that we can make the strong statement that the expectation is incorrect.  I've softened it to "not necessarily correct".  I can imagine scenarios where we wouldn't get the bimodal results in the presence of high formality.  Does this sound right to you?
% This variance paradox means that we need to reconsider the relationship between formalization and LLM output. The expectation that formalizing a task would make LLM output more predictable, is not necessarily correct. Rather, formalization results in higher mean quality, but can widen the distribution (higher variance). 
% % IS "CEILING" A common term IN STATs?
% Organizations using LLMs on 
% %schema-driven 
% modeling tasks should therefore expect, and design for, bimodal outcomes rather than gradual degradation.

% Thus, the study's most consequential finding is the refutation of H2: output variance \emph{decreases} monotonically as formalization decreases ($p < 10^{-76}$, Table~\ref{tab:levene}). This is the opposite of the degrees-of-freedom intuition. 
% %The mechanism is straightforward. 
% High-formal SQL produces bimodal outcomes. Either the model generates a query close to the gold standard or it diverges substantially (wrong joins, different strategies), creating high variance.
% %in overlap metrics.
% For low-formal management decision making, 
% %consistently cover the same conceptual territory (stakeholders, tradeoffs, risks)
% % Not SURE WHAT THE CONSISTency IS ABOUT OR WHERE it comes from BUT I PROBABLY JUST MISSED IT. 
% even when wording differs across runs, there was low variance.
% % SOMEthing funnY HERe.
% Formal tasks therefore yield higher mean quality \emph{and} higher variance; informal tasks yield lower mean quality \emph{but} more predictable outcomes. When using  LLMs for modeling tasks, users should expect, and design for, bimodal outcomes, by building in verification and retry mechanisms rather than assuming that high average quality implies consistent quality.
% % IS 'retry mechanisms' a term? or a Claude term?

The most consequential finding is the refutation of H2: output variance \textit{decreases} rather than increases with decreasing formalization. High-formal SQL tasks produced bimodal outcomes, where outputs were either close to the gold standard or substantially incorrect, whereas low-formal management tasks produced more consistently moderate outputs even when wording differed. Formalization therefore raises expected quality while widening the variance distribution, implying that schema-driven LLM applications require verification and retry mechanisms rather than assumptions of stable correctness.

% %\subsection{H1 Supported: The Formalization Gradient Is Real and Measurable}
% Hypothesis H1 was supported. The formalization level significantly predicts output quality, with all three overlap metrics declining monotonically from high to low ($p < 10^{-256}$, Table~\ref{tab:kruskal}).  The within-domain analysis (Section~\ref{sec:within-domain}) shows the same pattern across SQL difficulty tags and management complexity levels, confirming that formalization, not domain, drives the gradient. This empirical result extends the conceptual modeling prediction that semantic precision determines the quality of automated reasoning %\cite{DBLP:journals/isj/WandW93,olive2007conceptual}
% from logic-based to statistical reasoning systems, despite their fundamentally different inference mechanism. The gold-standard robustness analysis (Section~\ref{sec:gold-robustness}) rules out further differences.
% %stylistic contamination;
% % I didn't know whether 'stylistic contamination' is something well known.
% %the planned expert evaluation (Section~\ref{sec:human-results}) will provide independent validation using the unified 0--3 rubric.

H1 was supported: output quality declined sharply as formalization decreased, with the largest drop at the high-to-semi transition. Because the same pattern appeared both across domains and within domains of differing complexity, the results suggest that formalization is associated with LLM reasoning quality beyond domain alone (we hedge the causal reading in Section~\ref{sec:limitations}). This extends the conceptual modeling intuition that semantic precision improves automated reasoning from LLMs.


%\subsection{H3 Partially Supported: Response Length Mediates Quality Scores}

% The evidence for H3 is directionally consistent: Llama produces $0.6$--$8.4\times$ the gold-standard length vs.\ Mistral's $0.6$--$1.0\times$, with the quality gap between models being the largest where the length differential is largest (high-formal: $8.4\times$ vs.\ $1.0\times$) and vanishes where both produce equally short responses (low-formal: $0.6\times$ for both).
% % I THINK THE ABOVE SEnatence NEEDS A LITTLE REWRITING.
% H3 cannot be fully confirmed, however, without precision-based metrics that penalize irrelevant content. Future work should compute precision alongside the current recall-oriented metrics to quantify the mediating role of response length.

H3 received support. Longer responses were associated with higher recall-oriented scores, particularly for Llama at the high-formal level. Computing precision and $F_1$ alongside recall (Table~\ref{tab:precision}) resolves whether this verbosity premium reflects genuine understanding or lexical overlap inflation: it is inflation. At the high-formal level Mistral's concise outputs achieve substantially higher precision ($.823$ vs.\ $.547$) and $F_1$ ($.777$ vs.\ $.615$) than Llama's verbose ones, inverting the recall-based ranking. Response length therefore mediates the recall-based quality scores without a corresponding gain in precision, exactly as H3 predicts.

\subsection{Pairwise Comparison Anomaly and Consistency}

The pairwise results (Table~\ref{tab:pairwise}) reveal a nuanced pattern. The high$\,\to\,$semi and high$\,\to\,$low transitions are overwhelmingly significant with near-maximal effect sizes ($|r| > .94$) for every metric, so the dominant quality drop occurs at the high$\,\to\,$semi boundary. The semi$\,\to\,$low transition, by contrast, is small in magnitude for all metrics ($|\Delta| \leq .05$) but---owing to the large sample---statistically significant, and for word and key-term overlap it is a mild \emph{reversal}: low-formal tasks score slightly \emph{higher} than semi-formal ($\Delta = -.014$, $r = +.308$ for word; $\Delta = -.052$, $r = +.562$ for key-term). This likely reflects vocabulary differences (legal gold-standard answers use specialized terminology that models rarely reproduce, while management gold-standard answers use general analytical vocabulary that models match more easily) and shows the need for evaluation instruments that are independent of the vocabulary found in a domain.
%invariant 
%to domain vocabulary, 
This suggests the use of an assessment by an expert and within-domain analysis.

%\subsection{Consistency Remains an Open Problem}

The almost complete absence of output consistency across $K\!=\!5$ runs, even on objectively correct SQL tasks, highlights a fundamental challenge. Model temperature influences creativity (randomness) in its output. At $T\!=\!0.7$ (a typical general-purpose temperature setting), both models generate semantically similar, but lexically distinct outputs on every run. The 7\% perfect consistency achieved by Mistral on SQL (vs.\ 0\% for Llama) may reflect Mistral's more concise output style, which leaves less room for lexical variation.


%\subsection{Implications for Conceptual Modeling and Model-Driven Engineering}
% IS IT OK TO JUMPT TO MODEL-DRIVEN ENGINEERING? 

\subsection{Implications for Modeling}

%\paragraph{The automation boundary.}
High-formal tasks grounded in a conceptual model that led to a precise relational schema yield the highest quality, supporting LLM use for schema-driven code generation, query synthesis, and model transformation.
%\cite{olive2007conceptual}.
Semi-formal tasks (legal clause interpretation) show a substantial quality decline (word overlap drops by over 60\% relative to high-formal), indicating that LLM-assisted reasoning over semi-formal notations (UML constraints, business rules, structured requirements) is feasible, but requires careful verification. The quality plateau between semi-formal and low-formal levels suggests that
%dominant automation boundary sits 
the most important reliability boundary for LLM-based automation lies at the high$\,\to\,$semi transition, not the semi$\,\to\,$low transition.
%as might be expected.

%\paragraph{Formalization as a measurable investment.}
For formalization, the progression from low to high formalization results in a statistically significant gain in LLM output quality. This provides an empirical rationale for the long-standing conceptual modeling prescription to formalize where possible: investing in more formal representations (e.g., structured schemas, constraint languages for business rules) is effective when LLMs are used as reasoning engines. Combined with the variance paradox, the implication is that formal representations both raise expected quality and make verification \emph{necessary}, not merely desirable.
%\cite{DBLP:journals/isj/WandW93}.

We use \textit{context drift} to denote the tendency of generated outputs to diverge from the concepts, constraints, stakeholder relationships, or assumptions grounded in the input text. In random checking of the LLM outputs, we observed that context drift was minimal in the high-formal SQL tasks because the conceptual schema and executable semantics tightly constrained the solution space. In the semi-formal and low-formal tasks, drift occurred when LLMs generalized beyond the input text or introduced unsupported interpretations.

We use \textit{normative ambiguity} to denote situations in which multiple outputs are defensible because evaluation depends on competing stakeholder values, organizational priorities, or socially contingent judgments rather than formally decidable correctness criteria. Normative ambiguity is largely absent in the SQL tasks, bounded in the legal tasks by contractual language and precedent, and most pronounced in the management tasks where tradeoffs among stakeholders often lack a single objectively correct resolution.

These observations help explain the formalization gradient. As formalization decreases, the conceptual boundaries constraining interpretation weaken, increasing susceptibility to context drift and normative ambiguity even when outputs remain superficially coherent.


% ══════════════════════════════════════════════════════════════
\subsection{Limitations}
\label{sec:limitations}

The main results rely on text-overlap metrics as proxies for semantic quality; expert annotation under the unified $0$--$3$ rubric has been started but is not complete, so the findings lack independent human validation. The primary metrics are recall-oriented (dividing by $|W_{\text{gold}}|$ and not penalizing irrelevant content), which biases them toward verbose outputs; we mitigate this by reporting precision and $F_1$ alongside recall (Table~\ref{tab:precision}), which reveal that Llama's high-formal recall lead is verbosity-driven and reverses under precision. Overlap metrics of either kind remain coarser than execution-based or expert judgment, which the planned SQL execution evaluation and rubric annotation will supply.

Both evaluated models are $7$--$8$B parameters, and all runs use $T\!=\!0.7$. Generalization to $70$B+ models, where better instruction following and richer internal representations could alter the formalization gradient, is not warranted. Likewise, the variance paradox as reported here is strictly a $T\!=\!0.7$ phenomenon; a $T \in \{0, 0.3, 1.0\}$ sweep is needed to establish whether it persists across sampling regimes.

Semi-formal and low-formal gold standards were generated by Claude Sonnet 4.6 rather than by domain experts; expert validation is in progress, and until then the corresponding overlap metrics should be read as agreement with an LLM reference rather than as absolute correctness. Finally, the study is observational and cannot determine whether formalization \emph{causes} the observed quality differences or whether confounded task properties (domain, length, complexity) drive the gradient.

% ══════════════════════════════════════════════════════════════
\section{Conclusion and Future Work}
\label{sec:conclusion}

This research investigated how task formalization affects LLM performance in conceptual modeling applications. Across high-formal SQL generation, semi-formal legal clause interpretation, and low-formal management decision-making, output quality improved as formalization increased; output variance, however, followed the opposite pattern. High-formal tasks produced higher variance, while low-formal tasks produced lower but more consistently moderate scores. Formalization therefore raises expected quality but does not make LLM outputs more predictable; less formalized tasks are more susceptible to context drift and normative ambiguity because their conceptual boundaries and evaluation criteria are less explicit. For conceptual modeling, formal representations can improve LLM-assisted reasoning provided that verification mechanisms accompany them.

Several extensions follow directly. The dominant quality drop occurs at the high$\,\to\,$semi boundary, and both word and key-term recall show a mild reversal at the semi$\,\to\,$low transition; whether this reflects genuine non-monotonicity in the gradient or an artefact of domain vocabulary warrants investigation with richer instruments, including the expert $0$--$3$ rubric and within-domain complexity analysis. We have also prepared an obituary dataset (100 obituaries, two per US state, with annotated person/relation extractions corresponding to a conceptual model) as a semi-structured complement to the SQL benchmark; this dataset will join future experiments alongside additional datasets, larger models, and temperature sweeps.\footnote{Both datasets are available upon request.}

% ══════════════════════════════════════════════════════════════
\bibliographystyle{splncs04}
\bibliography{bibliography}

\end{document}